\section{Phụ lục: Cấu trúc mã nguồn và hướng dẫn tái tạo thực nghiệm}

Phụ lục này trình bày cấu trúc thư mục mã nguồn, vai trò của từng thành phần và hướng dẫn chạy lại các thực nghiệm trong báo cáo. Toàn bộ phần thực nghiệm được tổ chức trong thư mục \texttt{code/}.

\subsection{Cấu trúc thư mục}

Cấu trúc thư mục chính của project như sau:

\begin{lstlisting}
code/
|-- experiments/
|-- results/
|-- src/
|-- tests/
|-- README.md
|-- requirements.txt
|-- run_experiments.py
|-- run_main.py
\end{lstlisting}

Trong đó:

\begin{itemize}
    \item \texttt{experiments/}: chứa các script hoặc cấu hình phục vụ cho các thí nghiệm, chẳng hạn thí nghiệm chính, thí nghiệm thay đổi mức nhiễu và thí nghiệm thay đổi kích thước mẫu.
    
    \item \texttt{results/}: lưu kết quả đầu ra của quá trình thực nghiệm, bao gồm các bảng kết quả, file CSV và hình minh hoạ được sử dụng trong báo cáo.
    
    \item \texttt{src/}: chứa mã nguồn chính của chương trình, chẳng hạn các hàm sinh dữ liệu, mô hình ranking tuyến tính, hàm mất mát, hàm huấn luyện và các độ đo đánh giá.
    
    \item \texttt{tests/}: chứa các kiểm thử đơn giản nhằm đảm bảo các thành phần quan trọng như sinh cặp ưu tiên, tính loss và tính metric hoạt động đúng.
    
    \item \texttt{README.md}: mô tả ngắn gọn project, cách cài đặt môi trường và cách chạy chương trình.
    
    \item \texttt{requirements.txt}: liệt kê các thư viện Python cần thiết để chạy lại thực nghiệm.
    
    \item \texttt{run\_main.py}: script chạy thực nghiệm chính với thiết lập mặc định.
    
    \item \texttt{run\_experiments.py}: script chạy toàn bộ các thực nghiệm, bao gồm thí nghiệm chính, thí nghiệm theo mức nhiễu và thí nghiệm theo kích thước mẫu.
\end{itemize}

\subsection{Môi trường thực nghiệm}

Các thực nghiệm được cài đặt bằng Python. Các thư viện chính được sử dụng gồm:

\begin{itemize}
    \item \texttt{numpy}: dùng cho các phép toán vector, ma trận, sinh dữ liệu và cập nhật trọng số.
    \item \texttt{matplotlib}: dùng để vẽ các hình minh hoạ như đường cong loss, so sánh ranking và các biểu đồ thí nghiệm.
    \item \texttt{csv}: dùng để lưu kết quả thực nghiệm dưới dạng bảng.
    \item \texttt{pathlib}: dùng để quản lý đường dẫn thư mục và file.
\end{itemize}

Để cài đặt các thư viện cần thiết, chạy lệnh:

\begin{lstlisting}[language=bash]
pip install -r requirements.txt
\end{lstlisting}

\subsection{Cách chạy thực nghiệm chính}

Để chạy thực nghiệm chính với cấu hình mặc định, sử dụng lệnh:

\begin{lstlisting}[language=bash]
python run_main.py
\end{lstlisting}

Script này thực hiện các bước sau:

\begin{enumerate}[label=\textbf{Bước \arabic*.}, leftmargin=*]
    \item Sinh dữ liệu tổng hợp theo mô hình tuyến tính.
    \item Sinh các cặp ưu tiên từ điểm xếp hạng thật.
    \item Khởi tạo mô hình scoring tuyến tính.
    \item Huấn luyện mô hình bằng pairwise hinge loss.
    \item Tính các metric đánh giá như pairwise accuracy, ranking loss, Kendall tau distance và top-\(k\) overlap.
    \item Lưu kết quả và hình minh hoạ vào thư mục \texttt{results/}.
\end{enumerate}

\subsection{Cách chạy toàn bộ thực nghiệm}

Để chạy toàn bộ các thí nghiệm được trình bày trong báo cáo, sử dụng lệnh:

\begin{lstlisting}[language=bash]
python run_experiments.py
\end{lstlisting}

Script này bao gồm ba nhóm thực nghiệm:

\begin{itemize}
    \item \textbf{Thực nghiệm chính}: huấn luyện mô hình với thiết lập mặc định.
    \item \textbf{Thí nghiệm theo mức nhiễu}: thay đổi độ lệch chuẩn của nhiễu Gaussian để kiểm tra ảnh hưởng của nhiễu đến chất lượng ranking.
    \item \textbf{Thí nghiệm theo kích thước mẫu}: thay đổi số lượng mẫu huấn luyện để kiểm tra ảnh hưởng của kích thước dữ liệu.
\end{itemize}

Các kết quả sau khi chạy được lưu trong thư mục \texttt{results/}. Nếu các hình trong báo cáo được đặt trong thư mục \texttt{figures/}, có thể sao chép hoặc xuất các hình từ \texttt{results/} sang \texttt{figures/} để sử dụng trong file \LaTeX.

\subsection{Các file kết quả}

Sau khi chạy thực nghiệm, chương trình tạo ra các kết quả phục vụ cho phần báo cáo. Các file thường bao gồm:

\begin{itemize}
    \item Bảng kết quả thực nghiệm chính.
    \item Bảng kết quả khi thay đổi mức nhiễu.
    \item Bảng kết quả khi thay đổi kích thước mẫu.
    \item Hình đường cong pairwise hinge loss.
    \item Hình so sánh ranking thật và ranking dự đoán.
    \item Hình ảnh hưởng của nhiễu đến chất lượng ranking.
    \item Hình ảnh hưởng của kích thước mẫu đến chất lượng ranking.
\end{itemize}

Các bảng kết quả được dùng để tạo các bảng trong Section~9, còn các hình được dùng cho các minh hoạ như đường cong hội tụ và phân tích ảnh hưởng của tham số.

\subsection{Khả năng tái tạo}

Để đảm bảo kết quả có thể tái tạo, chương trình sử dụng random seed cố định trong quá trình sinh dữ liệu và khởi tạo mô hình. Các tham số chính của thực nghiệm, chẳng hạn số mẫu, số đặc trưng, learning rate, số epoch và hệ số regularization, được đặt rõ trong script thực nghiệm.

Khi chạy lại trên cùng môi trường và cùng random seed, các kết quả thu được sẽ gần như trùng khớp với kết quả đã trình bày trong báo cáo. Sai khác nhỏ có thể xuất hiện nếu phiên bản thư viện hoặc môi trường chạy khác nhau.

\subsection{Ghi chú về kiểm thử}

Thư mục \texttt{tests/} được dùng để lưu các kiểm thử cơ bản cho mã nguồn. Các kiểm thử này giúp đảm bảo rằng:

\begin{itemize}
    \item hàm sinh cặp ưu tiên trả về đúng số lượng cặp;
    \item nhãn cặp ưu tiên phù hợp với điểm xếp hạng thật;
    \item pairwise hinge loss không âm;
    \item ranking loss và pairwise accuracy có quan hệ đúng;
    \item các metric trả về giá trị trong miền hợp lệ.
\end{itemize}

Nếu có cấu hình kiểm thử bằng \texttt{pytest}, có thể chạy:

\begin{lstlisting}[language=bash]
pytest tests/
\end{lstlisting}

Nếu project không dùng \texttt{pytest}, phần kiểm thử có thể được chạy trực tiếp bằng các script Python trong thư mục \texttt{tests/}.